\documentclass[../Main.tex]{subfiles}

\begin{document}
%TODO: How are we defining a map??
Consider:
\begin{equation*}
    \vec{x_{n+1}} + \vec{F}(\vec{x_n})
\end{equation*}
where $\vec{F} : \R^n \mapsto \R^n$.

An obvious example of this is the discretisation of ODEs, if $\dvec{x} = \vec f(\vec x)$ then we can discretise as $\vec{X_{n+1}} = \vec{X_n} + \delta t f(\vec{x_n})$. Such maps can also occur in systems with discrete timing (this occurs in field like electronics and economics). We can also consider projections in high-dimensional space, where $\vec{X_n}$ are the points at which a high-dimensional system crosses a lower-dimensional plane.
\section{Motivational Examples of Maps}
\begin{example}[Logisic Map]
    \label{exLogisticMap}
    Consider the map:
    \begin{equation}
        x_{n+1} = \mu x_n(1-x_n)
        \label{eqnLogisticMap}
    \end{equation}
    which gives a parabola in the $x_n$-$x_{n+1}$ plane. See figure~\ref{figLogisticMap}. The way we read such a diagram is to start at a point $x_1$, move right to the line $x_{n+1} = x_n$, then move up to find $x_2$. Continuing the process, we see that there is an interesting fixed point where the straight line crosses the parabola.
\end{example}
\begin{figure}
    \centering
    \begin{tikzpicture}[scale=1]
        
    \end{tikzpicture}
    \caption{Diagram of the Logistic Map}
    \label{figLogisticMap}
\end{figure}
\begin{example}[Tent Map]
    \label{exTentMap}
    Consider the map:
    \begin{equation}
        x_{n+1} =
        \begin{cases}
            \mu x_n & 0 \leq x_n \leq \frac12 \\
            \mu(1-x_n) & \frac12 < x \leq 1
        \end{cases}
        \label{eqnTentMap}
    \end{equation}
    We see that this map has a fixed point for $\mu > 1$.
\end{example}
\begin{example}[Rotation Map]
    \label{exRotMap}
    Consider a map:
    \begin{equation*}
        x_{n+1} = x_n + \omega \mod 1
    \end{equation*}
    We see that if $\omega$ is rational then we get a periodic solution (with period the denominator). If $\omega$ is irrational then we do not get a periodic solution.
\end{example}
\begin{example}[Sawtooth (Bernoulli shift) map]
    \label{exSawToothMap}
    Consider a map:
    \begin{equation*}
        x_{n+1} = 2x_n \mod 1
    \end{equation*}
    Then this has a fixed point at $(1, 1)$. If we write $x_n$ in binary,
    \begin{equation*}
        x_n = 0.a_1a_2a_3\cdots
    \end{equation*}
    where $a_i \in \{0, 1\}$ are the binary digits after the point. Applying the map once will shift the digits as follows:
    \begin{equation*}
        x_{n+1} = 0.a_2a_3a_4\cdots
    \end{equation*}
    Solutions can therefore:
    \begin{itemize}
        \item hit zero eventually if $x_0$ has a terminating binary expansion;
        \item be periodic eventually if $x_0$ has a recurring binary expansion;
        \item be aperiodic if $x_0$ is irrational.
    \end{itemize}
\end{example}
\begin{example}[Baker Map]
    Consider the two-dimensional map:
    \begin{align*}
        x_n &= 2x_n \mod 1 \\
        y_n &= \begin{cases}
            \frac12 y_n & 0 \leq x_n < \frac12 \\
            \frac12 (y_n + 1) & \frac12 \leq x_n \leq 1
        \end{cases}
    \end{align*}
    %TODO: Draw the stretch/stack diagram.
\end{example}
\section{Stability of Maps}
\subsection{Definitions}
For this subsection, assume that $\vec{x_{n+1}} = \vec{F}(\vec{x_n})$ is a map. We will define our familiar concepts from ODEs in terms of maps.
\begin{definition}{Fixed point}
    A \underline{fixed point} is a value $\vec{x_0}$ such that $\vec{F}(\vec{x_0}) = \vec{x_0}$.
\end{definition}
\begin{definition}{Periodic orbit}
    A map has a \underline{periodic orbit} with \underline{period} $N$ if, for some $\vec{x_0}$, $\vec{F}^N(\vec{x_0}) = \vec{x_0}$ but $\vec F^k(\vec{x_0}) \neq \vec{x_0}$ for $k < N$.
\end{definition}
\begin{definition}{$N$-cycle}
    An \underline{$N$-cycle} $\{\vec{x_0}, \vec{x_1}, \cdots, \vec{x_{N-1}}\}$ is a set of points where $\vec{F}(\vec{x_i}) = \vec{x_{i+1}}$ and $\vec{F}(\vec{x_{N-1}}) = \vec{x_0}$. That is, it is all the points visited by a periodic orbit of period $N$.
\end{definition}
\begin{definition}{Invariant set}
    An \underline{invariant set} $\Lambda$ is a set with the property $\vec x \in \Lambda \implies \vec F(\vec x) \in \Lambda$. We can write $\vec F(\Lambda) = \Lambda$.
\end{definition}
\begin{definition}{Forward orbit}
    The \underline{forward orbit} $O^+(\vec x)$ is the set:
    \begin{equation*}
        O^+(\vec x) = \subsetselect{\vec F^n(\vec x)}{n \in \N}
    \end{equation*}
    and it is the set of all points visited by applying $\vec F$ to $\vec x$.
\end{definition}
The backward orbit is defined similarly.

We can now define our usual notions of stability -- the definitions of Lyapunov stability, quasi-asymptotic stability and asymptotic stability carry over directly.

\begin{definition}{Attractor}
    An \underline{attractor} is an invariant set which is asymptotically stable and has no nontrivial subset which is asymptotically stable.
\end{definition}
\subsection{Fixed Points}
Consider a fixed point $\vec{x_0}$ so $\vec{F}(\vec{x_0}) = \vec{x_0}$. We can linearise about this point:
\begin{equation*}
    \vec{x_n} - \vec{x_0} = \vec{y_n}
\end{equation*}
Then $\vec{x_{n+1}} = \vec{F}(\vec{x_n}) \implies \vec{y_{n+1}} + \vec{x_0} = \vec{F}(\vec{x_0} + \vec{y_n})$. Equating:
\begin{equation*}
    \vec{y_{n+1}} + \vec{x_0} = \vec F(\vec{x_0}) + A \vec{y_n} + O(\abs{\vec{y_n}}^2)
\end{equation*}
and so we have $\vec{y_{n+1}} = A\vec{y_n}$ where:
\begin{equation*}
    A_{ij} = \frac{\partial F_i}{\partial x_j}.
\end{equation*}
We can now characterise fixed points,
\begin{itemize}
    \item the fixed point is asymptotically stable if all eigenvalues of $A$ are less than 1 in magnitude;
    \item the fixed point is unstable if there is an eigenvalue of $A$ greater than 1 in magnitude;
    \item the fixed point is hyperbolic if none of the eigenvalues are $\pm 1$.
\end{itemize}
For one-dimensional maps, the matrix reduces to $A = \left.\frac{\partial F}{\partial x}\right|_{x = x_0}$, and so the dynamics are determined by the derivative of $F$. By considering cobweb diagrams,
\begin{itemize}
    \item $A > 1$ gives an unstable fixed point, with monotonic divergence;
    \item $A \in (0, 1)$ gives a stable fixed point, with monotonic convergence;
    \item $A \in (-1, 0)$ gives a stable fixed point, with oscillatory convergence (switching sides of $x_0$ each iteration);
    \item $A < -1$ gives an unstable fixed point, with oscillatory divergence.
\end{itemize}
We can see a clear similarity with nodes (monotonic convergence/divergence) and foci (oscillatory convergence/divergence).
\subsection{Stability of \texorpdfstring{$N$}{N}-Cycles}
Let $\vec{x_0}$ have a period $N$, and let its $N$-cycle be $\{\vec{x_i}\}_{i = 0}^{N-1}$. Let $\vec G(\vec x) = F^N(\vec x)$, so $\vec{x_0}$ is a fixed point of the map involving $\vec G$.

In one dimension,
\begin{align*}
    G'(x_0) &= \left.\frac{d}{dx} F^n(x)\right|_{x_0} \\
    &= F'(F^{N-1}(x_0)) F'(F^{N-2}(x_0)) \cdots F'(F(x_0)) F'(x_0) \\
    &= F'(x_{N-1}) F'(x_{N-2}) \cdots F'(x_0)
\end{align*}
and this will give the same expression for any of the $x_i$. This means stability of an $N$-cycle does not depend on where in the $N$-cycle we are.
\section{Bifurcations of Maps}
Consider a map with a parameter:
\begin{equation*}
    \vec{x_{n+1}} = \vec{F}(\vec{x_n}; \mu)
\end{equation*}
Let $\vec{x_0}$ be a fixed point of this map. Let $\lambda_{ij} = \left.\frac{\partial F_i}{\partial x_j}\right|_{\vec{x_0}}$. A bifurcation then happens when an eigenvalue of $\lambda$ is $\pm 1$ for $\mu = \mu_C$. We will take $\vec{x_0} = \vec0$ and $\mu_C = 0$ by perhaps shifting our coordinates and the parameter $\mu$.

In one dimension, we can have $\lambda = \pm 1$. We find that $F(0;0) = 0$ and $F'(0;0) = \pm 1$.
\subsection{One-Dimensional Bifurcations with Positive Eigenvalue}
In the case $\lambda = 1$, working in the vicinity of the bifurcation point, we can expand:
\begin{align*}
    x_{n+1} &=F(x_n;\mu) \\
    &= F(0;0) + F_x(0;0)x_n + F_\mu(0;0)\mu + \frac12 F_{xx}(0;0) x_n^2 \\
    &\qquad+ F_{x\mu}(0;0) x\mu + \frac12 F_{\mu\mu}(0;0)\mu^2 + \frac16 F_{xxx}(0;0) x_n^3 + \cdots \\
    &=x_n + F_\mu\mu + \frac12 F_{xx} x_n^2 + F_{x\mu}\mu x_n + \frac12 F_{\mu\mu}\mu^2 + \frac16 f_{xxx}x_n^3 + \cdots
\end{align*}
where in the last line we drop the $(0;0)$ and assume it for any terms involving $F$.

In the case that all coefficients of the expansion are non-vanishing, we keep the lowest order:
\begin{equation*}
    x_{n+1} = x_n + F_\mu \mu + \frac12 F_{xx} x_n^2
\end{equation*}
which has the normal form $x_{n+1} = x_n + \mu - x+n^2$, which looks like the saddle-node normal form for ODEs. Indeed, we see that the bifurcation diagram ($x_n$ vs $\mu$) is identical. The fixed points are $\pm \sqrt{\mu}$ with $x_n = \sqrt{\mu}$ stable and $-\sqrt{\mu}$ unstable.

In the case that $F(0;\mu) = 0$ for all $\mu$, collecting the lowest-order terms gives:
\begin{equation*}
    x_{n+1} = x_n + \frac12 F_{xx} x_n^2 + F_{x\mu} x\mu + O(3).
\end{equation*}
If $F_{xx} < 0, F_{x\mu} > 0$, then we can rescale to the normal form $x_{n+1} = x_n - x_n(x_n -\mu)$. We can identify this as a transcritical bifurcation: two fixed points exchange stability at $\mu = 0$.

Suppose now that $F_\mu = 0$ and $F_{xx} = 0$. Often this is due to symmetries of the equations such as $F(-x;\mu) = -F(x;\mu)$. Then the lowest-order terms in the expansion are:
\begin{equation*}
    x_{n+1} = x_n + F_{\mu x} \mu x_n + \frac16 F_{xxx}x_n^3 + O(\mu^2).
\end{equation*}
Then fixed points occur at $x_n = \pm \sqrt{-\frac{6f_{x\mu}}{f_{xxx}}\mu}$ and $x_n = 0$. The equation has the normal form:
\begin{equation*}
    x_{n+1} = x_n + x_n(\mu - ax_n)^2
\end{equation*}
with $a = \pm 1$. This corresponds to a pitchfork bifurcation.
\subsection{One-Dimensional Bifurcations with Negative Eigenvalue}
For $\lambda = -1$, we get some different behaviour. Assume that $F(0;\mu) = 0$, and now $F_x(0;0) = -1$.

In the previous subsection we found it was appropriate to use the scaling $x_n \sim \mu^{\frac12}$ to know which terms to keep.
\begin{equation*}
    x_{n+1} = -x_n + Bx_n^2 + A\mu x_n + Cx_n^3 + O(\mu^2)
\end{equation*}
where,
\begin{align*}
    B &= \frac12 F_{xx} & A &= F_{x\mu} & C &= \frac16 F_{xxx}.
\end{align*}
We look at $F^2$:
\begin{align*}
    x_{n+2} &= -x_{n+1} + Bx_{n+1}^2 + A\mu x_{n+1} + Cx_{n+1}^3 + O(\mu^2) \\
    &= x_n - Bx_n^2 - A\mu x_n - Cx_n^3 + B x_n^2 - 2B^2 x_n^3 \\
    &\qquad -A\mu x_n - Cx_n^3 + O(\mu^2) \\
    &= x_n - 2A\mu x_n - 2(B^2 + C)x_n^3 + O(\mu^2).
\end{align*}
Then the fixed points of $F^2$ occur when $x_n\left[A\mu + (B^2 + C)x_n^2\right] = 0$, so at $x_n = 0$ and at:
\begin{equation*}
    x_n =x_\pm = \pm \sqrt{-\frac{A}{B^2 + C}\mu}
\end{equation*}
Then the existence of fixed points of $F^2$ (but not of $F$) means that there are two-cycles corresponding to these values of $x_n$. This type of bifurcation is known as a \underline{period-doubling bifurcation}. We can calculate the stability of $x_n = 0$ for $F$:
\begin{equation*}
    F_x(0;0) = \left.(-1 + 2Bx_n + A\mu + 3Cx_n^2)\right|_{x = 0} = -1 + A\mu
\end{equation*}
That is, the point is stable for $A\mu > 0$ and unstable for $A\mu < 0$. For $x_n = x_\pm$, we can also find the stability (this time of $F^2$):
\begin{align*}
    F_x^2(0;0) &= 1 - 2A\mu - 6(B^2 + C)x_\pm^2 \\
    &= 1 - 2 A\mu - 6(B^2 + C) \left(-\frac{A\mu}{B^2 + C}\right) \\
    &= 1 + 4A\mu
\end{align*}
Then this has the opposite stability as $0$, $A\mu > 0$ gives an unstable periodic orbit, and $A\mu < 0$ gives an unstable periodic orbit.

Note that this period doubling bifurcation occurs for any setup with $\lambda < 0$.
\end{document}